\chapter{60}\label{section-60}

\section{Mailand (IT), Mittwoch, 22.
August}\label{mailand-it-mittwoch-22.-august-1}

«Vorne weg, alle drei wirkten sehr betroffen. Und alle drei meinten,
keine Vorstellung zu haben, was sich in der Nacht ereignet hat.»

Riva setzte ihre Lesebrille auf und sammelte sich. Sie sass zwischen
Sömnez und Rossi. Mantovani ihnen gegenüber.

«Nach Maria Greco müsse es kurz vor Mitternacht gewesen sein. Giorgia
Sala habe gesagt, sie müsse austreten. Dann sei sie umständlich
aufgestanden und torkelnd in der Dunkelheit verschwunden. Und sie habe
ihr noch nachgerufen, ob sie mitkommen solle, und diese habe
geantwortet, nein, sie gehe alleine. Sie fügte hinzu, sie erinnere sich,
wie heiser Sala dabei gelacht hatte. Es sei ihr aufgefallen, dass diese
sehr merkwürdig drauf gewesen sei. Sie wisse nicht genau, warum.
Getrunken habe sie schon genug, aber trotzdem. Und ja, Bastiano di
Martino sei auch da gewesen. Der sei zwei Mal mit einem Mädchen im
kleinen Wäldchen verschwunden. Aber das sei nichts Neues. Mit Giorgia
Sala zusammen hätte sie ihn nicht gesehen. Der habe gar keine Zeit dazu
gehabt. Vielleicht sei er kurz mit Sala in Berührung gekommen, als er
die mitgebrachten Getränke verteilte. Und er sei dann für seine
Verhältnisse früh gegangen. Ja, auch Louis Vincent sei nach längerer
Zeit wieder einmal auf den \emph{Roccia} gekommen. Giorgia Salas
Exfreund. Die hätten sich vor einigen Wochen getrennt. Darum habe er
wohl Abstand gebraucht, denke sie. Zu Anfang habe er neben Valentina
Ricci am Feuer gesessen. Auch er sei vor Mitternacht gegangen. Nein, von
Drogen wisse sie nichts. Alkohol ja, vielleicht der eine oder andere
Joint, aber anderes? Davon wisse sie wirklich nichts, und…»

Riva holte Luft.

«Franco Espositos Aussagen deckten sich mit Maria Grecos Antworten. Die
dritte, Valentina Ricci, vernahm ich als Letzte,» Riva legte die Stirn
in Falten, «bei der Begrüssung fiel mir ihre schweissnasse Hand auf. Ja,
und während des Gesprächs wich sie meinem Blick auffallend oft aus. Das
kann natürlich verschiedene Gründe haben. Sie scheint mir jedenfalls
eine schüchterne, leise Person zu sein.»

Mantovani zeigte auf Rossi.

«Welche der beiden Mädchen hast du auf Körpersprache beobachtet?»

«Valentina Ricci.»

«Gut. Deine Beobachtungen brauchen wir gleich als Ergänzung,» er wandte
sich wieder Riva zu und bedeutete ihr, weiter zu sprechen.

«Valentina Ricci rutschte auffallend auf dem Stuhl herum und legte ihre
Hände auf die Knie. Mir schien, sie wolle sich verstecken. Dann fragte
ich sie, wie gut sie Bastiano di Martino kenne. Sie erschrak sichtlich
und sprach ganz leise. Sie kenne ihn nicht sehr gut. Er sei ein Freund
von Giorgia Sala, seit sie Kinder seien. Sie sei ihm nicht nahe. Er sei
ein Frauenheld, ziemlich laut und einfach einer, der auch zur Clique
gehöre.»

Riva machte eine Pause.

«Hm,» sie richtete ihren Blick auf ihr aufgeschlagenes Notizbuch, «wie
bereits Maria Greco aussagte, habe Louis Vincent mit ihr am Feuer
gesessen. Gesprochen hätten sie kaum miteinander, weil die Musik sehr
laut war. Und auch sie sagte, Vincent sei vor Mitternacht gegangen.
Schliesslich fragte ich sie, ob in dieser Nacht ein Mann mit dabei
gewesen sei, ungefähr 40 Jahre alt, mit einer auffallend hohen Stimme,
etwa eins siebzig-fünfundsiebzig gross, und mit einem speziell
eindringlichen Blick. Ich konnte direkt zusehen, wie sie sich
entspannte. Ihre Stimme klang plötzlich sehr klar.»

Die drei Männer blickten Riva erwartungsvoll an.

«Nein, da sei keiner gewesen, der auf die Beschreibung passe, sagte sie.
Ihr seien alle bekannt gewesen, ausser zwei junge Mädchen, die di
Martino eingeladen hatte. Valentina Ricci hat dazu vehement den Kopf
geschüttelt, ja, und Maria Grecos und Franco Espositos Reaktionen auf
diese Frage waren ebenfalls deckungsgleich.»

Riva schaute in die Runde.

«Nach unserer Verabschiedung wirkte Valentina Ricci auffallend
erleichtert und verliess den Raum schnell. Mir war, als müsse sie sich
zurückhalten, nicht raus zu rennen.»

Mantovani nickte und gab Sömnez das Wort. Der fügte seine Beobachtungen
hinzu. Sie ergänzten Rivas Bild von Valentina Ricci. Sie waren sich
einig. Das Verhalten der jungen Frau zeigte mindestens Unsicherheit. Je
nach Verlauf der weiteren Ermittlungen würden sie sie erneut vorladen.

Mantovani rieb sich die Stirn und stand auf. Er kniff die Augen zusammen
und setzte sich wieder.

«Ich danke euch. Das grösste Rätsel ist noch immer der Mann, der den
Sturz meldete,» er starrte in die Ferne, «was ist der Grund seines
Auftritts? Warum meldet er den Unfall, oder was immer es war, hier
direkt auf der Questur? Warum getarnt? Was versteckt er? Was hat der
damit zu tun? Und wenn er’s wirklich nur gesehen hat, warum tritt er
nicht als Privatperson auf und zeigt sich.»

Rossi hob die Augenbrauen. Sömnez zuckte mit den Schultern und Emma Riva
bündelte ihre Unterlagen.

«Nun denn, wir kommen nicht drumherum,» Mantovani streckte den Rücken
durch, «bitte, Emma, veranlasse die Veröffentlichung der Videoaufnahme
des Mannes in den Medien. Vielleicht erkennt ihn jemand,» er wandte sich
direkt Sömnez zu, «Melik, du fragst bei allen drei nach, ob sie Näheres
über die Trennung von diesem Vincent und Giorgia Sala wissen, ob sie
einvernehmlich war, oder ein Drama, ob sie vielleicht wissen, wer wen
verlassen hat…du weisst schon, kannst es telefonisch machen. Für’s
Erste. Stelle dieselben Fragen auch den Eltern Sala.»

Sömnez nickte.
